% GENERATED FILE. Edit canonical lessons, then run npm run generate:books.
% canonical-source-hash: fnv1a64:ba622bb6228c37d6
% canonical-lessons: PA-C336-karigar, PA-C336-mazdur, PA-C336-kasai, PA-C336-daraivar, PA-C336-murammat, PA-R336-review-330-336, PA-R336-review-final-336

\chapter{Work in the Bazaar}
\label{ch:pa-a2-work-in-the-bazaar}
\hlchaptermodality{336}

\begin{chapteropening}
\textbf{By the end of this chapter:} \emph{I can say what work people do in the bazaar.}

Complete the last lesson of chapter 336: I can say what work people do in the bazaar.
\end{chapteropening}

\section[\texorpdfstring{\pa{ਕਾਰੀਗਰ} (kārīgar)}{ਕਾਰੀਗਰ (kārīgar)}]{\pa{ਕਾਰੀਗਰ} (kārīgar) — a craftsman}

\label{lesson:PA-C336-karigar}



\begin{quote}
Before the new one: say the Punjabi for to hammer into shape, then the Punjabi for a mechanic.
\end{quote}

\subsection*{You'll want to know: \pa{ਕਾਰੀਗਰ}}
\textbf{\pa{ਕਾਰੀਗਰ}} — \emph{kārīgar} — \textquotedblleft{}a craftsman\textquotedblright{}.

\begin{quote}
\pa{ਇਹ} \pa{ਕਾਰੀਗਰ} \pa{ਲੱਕੜ} \pa{ਦਾ} \pa{ਕੰਮ} \pa{ਕਰਦਾ} \pa{ਹੈ।}
\end{quote}

\begin{quote}
\emph{eh kārīgar lakkaṛ dā kamm kardā hai.}
\end{quote}

\begin{quote}
This craftsman does woodwork.
\end{quote}

A masculine noun.

\begin{cousinweb}
\textbf{\pa{ਕਾਰੀਗਰ}} (\emph{kārīgar}) is a loan, usually traced to Persian \emph{kārīgar}, \textquotedblleft{}a craftsman\textquotedblright{}.
\end{cousinweb}

\begin{grammarlens}[title={What you've built}]
One more word for this chapter.
\end{grammarlens}

\subsection*{Your turn}
\begin{itemize}
\raggedright
  \item \emph{Say it:} \emph{kārīgar}
  \item \emph{Say it:} \emph{kārīgar}, once more
  \item \emph{Say it:} \emph{eh kārīgar lakkaṛ dā kamm kardā hai.} — This craftsman does woodwork.
  \item \emph{Recall:} say the Punjabi for to hammer into shape, then the Punjabi for a mechanic, then say \emph{kārīgar} again
\end{itemize}

\begin{tcolorbox}[breakable,colback=teal!4,colframe=teal!35!black,title={Before you move on}]
What does \pa{ਕਾਰੀਗਰ} mean? (\textquotedblleft{}A craftsman\textquotedblright{}.) Say it once more.
\end{tcolorbox}
\section[\texorpdfstring{\pa{ਮਜ਼ਦੂਰ} (mazdūr)}{ਮਜ਼ਦੂਰ (mazdūr)}]{\pa{ਮਜ਼ਦੂਰ} (mazdūr) — a labourer}

\label{lesson:PA-C336-mazdur}



\begin{quote}
Before the new one: say the Punjabi for a mechanic, then the Punjabi for a craftsman.
\end{quote}

\subsection*{You'll want to know: \pa{ਮਜ਼ਦੂਰ}}
\textbf{\pa{ਮਜ਼ਦੂਰ}} — \emph{mazdūr} — \textquotedblleft{}a labourer\textquotedblright{}.

\begin{quote}
\pa{ਮਜ਼ਦੂਰ} \pa{ਸਵੇਰ} \pa{ਤੋਂ} \pa{ਸ਼ਾਮ} \pa{ਤੱਕ} \pa{ਇੱਟਾਂ} \pa{ਚੁੱਕਦੇ} \pa{ਹਨ।}
\end{quote}

\begin{quote}
\emph{mazdūr saver ton shām tak iṭṭāṁ chukkde han.}
\end{quote}

\begin{quote}
The labourers carry bricks from morning till evening.
\end{quote}

A masculine noun.

\begin{cousinweb}
\textbf{\pa{ਮਜ਼ਦੂਰ}} (\emph{mazdūr}) is a loan, usually traced to Persian \emph{mazdūr}, \textquotedblleft{}a hired labourer\textquotedblright{}.
\end{cousinweb}

\begin{grammarlens}[title={What you've built}]
One more word for this chapter.
\end{grammarlens}

\subsection*{Your turn}
\begin{itemize}
\raggedright
  \item \emph{Say it:} \emph{mazdūr}
  \item \emph{Say it:} \emph{mazdūr}, once more
  \item \emph{Say it:} \emph{mazdūr saver ton shām tak iṭṭāṁ chukkde han.} — The labourers carry bricks from morning till evening.
  \item \emph{Recall:} say the Punjabi for a mechanic, then the Punjabi for a craftsman, then say \emph{mazdūr} again
  \item \emph{Write it:} copy \textbf{\pa{ਮਜ਼ਦੂਰ}} once, checking the whole word
\end{itemize}

\begin{tcolorbox}[breakable,colback=teal!4,colframe=teal!35!black,title={Before you move on}]
What does \pa{ਮਜ਼ਦੂਰ} mean? (\textquotedblleft{}A labourer\textquotedblright{}.) Say it once more.
\end{tcolorbox}
\section[\texorpdfstring{\pa{ਕਸਾਈ} (kasāī)}{ਕਸਾਈ (kasāī)}]{\pa{ਕਸਾਈ} (kasāī) — a butcher}

\label{lesson:PA-C336-kasai}



\begin{quote}
Before the new one: say the Punjabi for a craftsman, then the Punjabi for a labourer.
\end{quote}

\subsection*{You'll want to know: \pa{ਕਸਾਈ}}
\textbf{\pa{ਕਸਾਈ}} — \emph{kasāī} — \textquotedblleft{}a butcher\textquotedblright{}.

\begin{quote}
\pa{ਕਸਾਈ} \pa{ਦੀ} \pa{ਦੁਕਾਨ} \pa{ਬਾਜ਼ਾਰ} \pa{ਦੇ} \pa{ਕੋਨੇ} \pa{ਉੱਤੇ} \pa{ਹੈ।}
\end{quote}

\begin{quote}
\emph{kasāī dī dukān bāzār de kone utte hai.}
\end{quote}

\begin{quote}
The butcher's shop is on the corner of the bazaar.
\end{quote}

A masculine noun.

\begin{cousinweb}
\textbf{\pa{ਕਸਾਈ}} (\emph{kasāī}) is a loan, usually traced through Persian to Arabic \emph{qaṣṣāb}, \textquotedblleft{}a butcher\textquotedblright{}.
\end{cousinweb}

\begin{grammarlens}[title={What you've built}]
One more word for this chapter.
\end{grammarlens}

\subsection*{Your turn}
\begin{itemize}
\raggedright
  \item \emph{Say it:} \emph{kasāī}
  \item \emph{Say it:} \emph{kasāī}, once more
  \item \emph{Say it:} \emph{kasāī dī dukān bāzār de kone utte hai.} — The butcher's shop is on the corner of the bazaar.
  \item \emph{Recall:} say the Punjabi for a craftsman, then the Punjabi for a labourer, then say \emph{kasāī} again
  \item \emph{Write it:} copy \textbf{\pa{ਕਸਾਈ}} once, checking the whole word
\end{itemize}

\begin{tcolorbox}[breakable,colback=teal!4,colframe=teal!35!black,title={Before you move on}]
What does \pa{ਕਸਾਈ} mean? (\textquotedblleft{}A butcher\textquotedblright{}.) Say it once more.
\end{tcolorbox}
\section[\texorpdfstring{\pa{ਡਰਾਈਵਰ} (ḍarāīvar)}{ਡਰਾਈਵਰ (ḍarāīvar)}]{\pa{ਡਰਾਈਵਰ} (ḍarāīvar) — a driver}

\label{lesson:PA-C336-daraivar}



\begin{quote}
Before the new one: say the Punjabi for a labourer, then the Punjabi for a butcher.
\end{quote}

\subsection*{You'll want to know: \pa{ਡਰਾਈਵਰ}}
\textbf{\pa{ਡਰਾਈਵਰ}} — \emph{ḍarāīvar} — \textquotedblleft{}a driver\textquotedblright{}.

\begin{quote}
\pa{ਸਾਡਾ} \pa{ਡਰਾਈਵਰ} \pa{ਬੱਸ} \pa{ਬਹੁਤ} \pa{ਹੌਲੀ} \pa{ਚਲਾਉਂਦਾ} \pa{ਹੈ।}
\end{quote}

\begin{quote}
\emph{sāḍā ḍarāīvar bass bahut haulī chalāundā hai.}
\end{quote}

\begin{quote}
Our driver drives the bus very slowly.
\end{quote}

A masculine or feminine noun, by the person meant.

\begin{cousinweb}
\textbf{\pa{ਡਰਾਈਵਰ}} (\emph{ḍarāīvar}) is borrowed from English \emph{driver}.
\end{cousinweb}

\begin{grammarlens}[title={What you've built}]
One more word for this chapter.
\end{grammarlens}

\subsection*{Your turn}
\begin{itemize}
\raggedright
  \item \emph{Say it:} \emph{ḍarāīvar}
  \item \emph{Say it:} \emph{ḍarāīvar}, once more
  \item \emph{Say it:} \emph{sāḍā ḍarāīvar bass bahut haulī chalāundā hai.} — Our driver drives the bus very slowly.
  \item \emph{Recall:} say the Punjabi for a labourer, then the Punjabi for a butcher, then say \emph{ḍarāīvar} again
  \item \emph{Write it:} copy \textbf{\pa{ਡਰਾਈਵਰ}} once, checking the whole word
\end{itemize}

\begin{tcolorbox}[breakable,colback=teal!4,colframe=teal!35!black,title={Before you move on}]
What does \pa{ਡਰਾਈਵਰ} mean? (\textquotedblleft{}A driver\textquotedblright{}.) Say it once more.
\end{tcolorbox}
\section[\texorpdfstring{\pa{ਮੁਰੰਮਤ} (murammat)}{ਮੁਰੰਮਤ (murammat)}]{\pa{ਮੁਰੰਮਤ} (murammat) — a repair}

\label{lesson:PA-C336-murammat}



\begin{quote}
Before the new one: say the Punjabi for a butcher, then the Punjabi for a driver.
\end{quote}

\subsection*{You'll want to know: \pa{ਮੁਰੰਮਤ}}
\textbf{\pa{ਮੁਰੰਮਤ}} — \emph{murammat} — \textquotedblleft{}a repair\textquotedblright{}.

\begin{quote}
\pa{ਪਿਛਲੇ} \pa{ਮਹੀਨੇ} \pa{ਅਸੀਂ} \pa{ਘਰ} \pa{ਦੀ} \pa{ਮੁਰੰਮਤ} \pa{ਕੀਤੀ।}
\end{quote}

\begin{quote}
\emph{pichhle mahīne asīṁ ghar dī murammat kītī.}
\end{quote}

\begin{quote}
Last month we repaired the house.
\end{quote}

A feminine noun.

\begin{cousinweb}
\textbf{\pa{ਮੁਰੰਮਤ}} (\emph{murammat}) is a loan, usually traced through Persian to Arabic \emph{marammat}, \textquotedblleft{}repair\textquotedblright{}.
\end{cousinweb}

\begin{grammarlens}[title={What you've built}]
One more word for this chapter.
\end{grammarlens}

\subsection*{Your turn}
\begin{itemize}
\raggedright
  \item \emph{Say it:} \emph{murammat}
  \item \emph{Say it:} \emph{murammat}, once more
  \item \emph{Say it:} \emph{pichhle mahīne asīṁ ghar dī murammat kītī.} — Last month we repaired the house.
  \item \emph{Recall:} say the Punjabi for a butcher, then the Punjabi for a driver, then say \emph{murammat} again
  \item \emph{Write it:} copy \textbf{\pa{ਮੁਰੰਮਤ}} once, checking the whole word
\end{itemize}

\begin{tcolorbox}[breakable,colback=teal!4,colframe=teal!35!black,title={Before you move on}]
What does \pa{ਮੁਰੰਮਤ} mean? (\textquotedblleft{}A repair\textquotedblright{}.) Say it once more.
\end{tcolorbox}
\section[(dialogue)]{Looking back: from Brothers' and Sisters' Spouses to Work in the Bazaar}

\label{lesson:PA-R336-review-330-336}



\begin{quote}
Say the Punjabi for a brother's wife and a sister's husband.
\end{quote}

\subsection*{Your turn}
Say these aloud:

\begin{itemize}
\raggedright
  \item a brother's wife, a sister's husband, a husband's sister, a husband's younger brother's wife, a husband's elder brother's wife
  \item a father's younger brother, a father's elder brother, a mother's brother, a mother's sister, a mother's sister's husband
  \item a brother's son, a sister's son, a neighbour, a woman's female friend, to call out to someone
  \item a barber, to shave, a washerman, a cobbler, a gardener
  \item wool, to spin thread, to knit, a weaver, to dye
  \item gold, a goldsmith, iron, to hammer into shape, a mechanic
  \item a craftsman, a labourer, a butcher, a driver, a repair
\end{itemize}

\begin{tcolorbox}[breakable,colback=teal!4,colframe=teal!35!black,title={Before you move on}]
A brother's wife? (\textbf{\pa{ਭਰਜਾਈ}}, \emph{bharjāī}.) A sister's husband? (\textbf{\pa{ਜੀਜਾ}}, \emph{jījā}.) A husband's sister? (\textbf{\pa{ਨਣਾਨ}}, \emph{naṇān}.) A husband's younger brother's wife? (\textbf{\pa{ਦਰਾਣੀ}}, \emph{darāṇī}.) A husband's elder brother's wife? (\textbf{\pa{ਜਠਾਣੀ}}, \emph{jaṭhāṇī}.) A father's younger brother? (\textbf{\pa{ਚਾਚਾ}}, \emph{chāchā}.) A father's elder brother? (\textbf{\pa{ਤਾਇਆ}}, \emph{tāiā}.) A mother's brother? (\textbf{\pa{ਮਾਮਾ}}, \emph{māmā}.) A mother's sister? (\textbf{\pa{ਮਾਸੀ}}, \emph{māsī}.) A mother's sister's husband? (\textbf{\pa{ਮਾਸੜ}}, \emph{māsaṛ}.) A brother's son? (\textbf{\pa{ਭਤੀਜਾ}}, \emph{bhatījā}.) A sister's son? (\textbf{\pa{ਭਾਣਜਾ}}, \emph{bhāṇjā}.) A neighbour? (\textbf{\pa{ਗੁਆਂਢੀ}}, \emph{guāṇḍhī}.) A woman's female friend? (\textbf{\pa{ਸਹੇਲੀ}}, \emph{sahelī}.) To call out to someone? (\textbf{\pa{ਪੁਕਾਰਨਾ}}, \emph{pukārnā}.) A barber? (\textbf{\pa{ਨਾਈ}}, \emph{nāī}.) To shave? (\textbf{\pa{ਮੁੰਨਣਾ}}, \emph{munnṇā}.) A washerman? (\textbf{\pa{ਧੋਬੀ}}, \emph{dhobī}.) A cobbler? (\textbf{\pa{ਮੋਚੀ}}, \emph{mochī}.) A gardener? (\textbf{\pa{ਮਾਲੀ}}, \emph{mālī}.) Wool? (\textbf{\pa{ਉੱਨ}}, \emph{unn}.) To spin thread? (\textbf{\pa{ਕੱਤਣਾ}}, \emph{kattṇā}.) To knit? (\textbf{\pa{ਬੁਣਨਾ}}, \emph{buṇnā}.) A weaver? (\textbf{\pa{ਜੁਲਾਹਾ}}, \emph{julāhā}.) To dye? (\textbf{\pa{ਰੰਗਣਾ}}, \emph{raṅgṇā}.) Gold? (\textbf{\pa{ਸੋਨਾ}}, \emph{sonā}.) A goldsmith? (\textbf{\pa{ਸੁਨਿਆਰ}}, \emph{suniār}.) Iron? (\textbf{\pa{ਲੋਹਾ}}, \emph{lohā}.) To hammer into shape? (\textbf{\pa{ਘੜਨਾ}}, \emph{ghaṛnā}.) A mechanic? (\textbf{\pa{ਮਿਸਤਰੀ}}, \emph{misatrī}.) A craftsman? (\textbf{\pa{ਕਾਰੀਗਰ}}, \emph{kārīgar}.) A labourer? (\textbf{\pa{ਮਜ਼ਦੂਰ}}, \emph{mazdūr}.) A butcher? (\textbf{\pa{ਕਸਾਈ}}, \emph{kasāī}.) A driver? (\textbf{\pa{ਡਰਾਈਵਰ}}, \emph{ḍarāīvar}.) A repair? (\textbf{\pa{ਮੁਰੰਮਤ}}, \emph{murammat}.)
\end{tcolorbox}
\section[(dialogue)]{Work in the Bazaar — once more}

\label{lesson:PA-R336-review-final-336}



\begin{quote}
Say the Punjabi for a craftsman and a labourer.
\end{quote}

\subsection*{Your turn}
\begin{itemize}
\raggedright
  \item \emph{Say it:} a craftsman, a labourer, a butcher, a driver, a repair
\end{itemize}

\begin{tcolorbox}[breakable,colback=teal!4,colframe=teal!35!black,title={Before you move on}]
A craftsman? (\textbf{\pa{ਕਾਰੀਗਰ}}, \emph{kārīgar}.) A labourer? (\textbf{\pa{ਮਜ਼ਦੂਰ}}, \emph{mazdūr}.) A butcher? (\textbf{\pa{ਕਸਾਈ}}, \emph{kasāī}.) A driver? (\textbf{\pa{ਡਰਾਈਵਰ}}, \emph{ḍarāīvar}.) A repair? (\textbf{\pa{ਮੁਰੰਮਤ}}, \emph{murammat}.)
\end{tcolorbox}
